\documentclass[12pt]{article}

\usepackage[margin=1in]{geometry}
\usepackage{amsmath,amssymb,amsthm}
\usepackage{natbib}
\usepackage{booktabs}
\usepackage{hyperref}
\usepackage{setspace}

\newtheorem{proposition}{Proposition}
\newtheorem{assumption}{Assumption}
\theoremstyle{remark}
\newtheorem{remark}{Remark}

\onehalfspacing

\title{From Data Products to Inference:\\ The Econometrics of Fuzzy Joins}
\author{Gaurav Sood}
\date{\today}

\begin{document}
\maketitle

\begin{abstract}
Data products built on fuzzy joins often ship a single canonical linkage while concealing the candidate-match graph from which that linkage was constructed. This note studies downstream regression when the analyst can observe the candidate graph. Two natural single-dataset resolutions fail for different reasons. Edge expansion introduces an algebraic variance penalty that mechanically shrinks the slope relative to a score-collapsed regression. Score averaging, by contrast, induces nonclassical measurement error whose bias need not attenuate and can have either sign. Which linkage errors matter depends on the estimand: joins that assign treatment status, joins that attach outcomes to pre-assigned groups, and level estimands face different error costs. I discuss two practical corrections---multiple imputation over candidate assignments and validation-sample regression calibration---and then show that both may miss a second problem: upstream contamination when the covariates shipped in the target table were themselves created by an earlier fuzzy join. The note's main message is that linkage uncertainty is not a single scalar property of a linked file; it is estimand-specific, and some of it is inherited from the data product itself.
\end{abstract}

\section{Introduction}

Empirical researchers increasingly work with linked data products assembled upstream from administrative or transactional files. In many applications the linking step is not exact: records are joined using names, dates, locations, and other quasi-identifiers that produce a set of plausible candidate matches rather than a single certainty. The analyst, however, usually receives a clean final file and proceeds as if each row were linked without error.

This note asks a downstream question. Suppose the data producer ships not only a canonical linkage but the candidate-match graph itself: for each source record, all plausible target records together with a pairwise score. How should an analyst convert that object into a regression dataset, and what happens if they do so naively?

The existing literature on regression with linked data is substantial. Early bias-correction work includes \citet{scheuren1993,scheuren1997} and \citet{lahiri2005}. Subsequent contributions address probability-linked data, incomplete linkage, probabilistic multi-linkage, correlated linkage error, and linked-data secondary analysis more broadly \citep{chambers2009,kim2012multi,kim2012incomplete,kim2015,han2019,chambers2020,wang2022,doidge2019}. Multiple-imputation and primary-analysis approaches include \citet{goldstein2012,harron2014,hof2012,hof2015,hof2017,tancredi2011,tancredi2015,gutman2013,steorts2016,dalzell2018}. More recent post-linkage work includes \citet{slawski2025}; \citet{slawski2021} is a closely related precursor for partially shuffled linear regression rather than a general linked-data framework. A companion paper \citep{sood2025selection} formalizes the join itself as a selection problem over feasible linked tables, connecting to the finite selection model literature; this note takes up the downstream inference question that follows.

My contribution is narrower. I focus on the econometric problem created when a data product ships a many-to-many candidate graph for a fuzzy geographic join. I make four points. First, the inferential consequences of linkage error depend on the estimand. Second, the two most obvious one-shot resolutions of the candidate graph produce distinct pathologies. Third, multiple imputation and validation-based calibration are useful but answer different questions. Fourth, even perfect treatment of the analyst's own candidate graph does not solve upstream linkage error frozen into the shipped covariates.

\section{Setup and notation}

Let $S=\{1,\ldots,n\}$ index source records and $T=\{1,\ldots,K\}$ index target records. For source unit $i$, let the true target be $v(i)\in T$. The regression of interest is
\begin{equation}
Y_i=\beta_0+\beta_1 X_{v(i)}+\eta_{v(i)}+\varepsilon_i,
\label{eq:dgp}
\end{equation}
where $X_k$ is a target-level covariate, $\eta_k$ is a target-level shock, and $\varepsilon_i$ is idiosyncratic noise.

A fuzzy matcher produces a candidate graph $\mathcal G=(S,T,E,s)$, where $E\subseteq S\times T$ is the set of candidate edges and $s:E\to[0,1]$ is a pairwise score. For source unit $i$, define the candidate set
\[
\mathcal C_i=\{(j,s_{ij}):(i,j)\in E\}.
\]
I distinguish throughout between the raw pairwise score $s_{ij}$ and a calibrated match probability
\[
p_{ij}\approx P(v(i)=j\mid \mathcal C_i),
\]
which is an additional object that must be estimated; it is not implied by the score itself.

A central theme of the paper is the mismatch between the \emph{edge-level} nature of the candidate graph and the \emph{row-level} structure of the regression dataset \citep[see][for a general treatment of this mismatch]{sood2025selection}. The graph contains many candidate edges per source unit. The regression requires one covariate and one cluster assignment per source unit.

\section{Which linkage errors matter depends on the estimand}

The first mistake analysts make is to talk about ``linkage error'' as though it had a universal inferential cost. It does not.

\paragraph{Treatment-assignment joins.}
If the linkage determines treatment status, a false match contaminates both groups: one unit is moved into the wrong group and another is displaced from the correct one. In that setting, false positives and false negatives both matter directly for bias.

\paragraph{Outcome-attachment joins within pre-assigned groups.}
If the groups are fixed before linkage and the join only attaches outcomes within group, then within-group permutations are benign for unweighted group means. Only errors that cross an estimand-relevant boundary matter for bias.

\paragraph{Level estimands.}
If the estimand is a level for a specific population, every wrong attachment matters.

The practical implication is that the ``best'' resolution of a candidate graph is estimand-specific. A treatment-assignment design typically requires a much harsher trade-off against false matches than a design in which ambiguity remains within groups.

When calibrated probabilities are available, the resolution problem can be written as a maximum-weight matching problem. Let the expected utility of selecting edge $(i,j)$ be
\[
a_{ij}=u_{\mathrm{TP}}\,p_{ij}-u_{\mathrm{FP}}(1-p_{ij}),
\]
where $u_{\mathrm{TP}}>0$ is the value of a true positive and $u_{\mathrm{FP}}>0$ is the cost of a false positive. Then the linkage rule solves
\begin{equation}
\max_{L\subseteq E} \sum_{(i,j)\in L} a_{ij}
\label{eq:decision}
\end{equation}
subject to the application's assignment constraints. Edges with $a_{ij}\le 0$ can be screened out ex ante, but the final solution remains a constrained optimization problem rather than a pure edgewise threshold rule.

\section{Two naive resolutions of the candidate graph}

\subsection{Expanded edge-level regression}

A natural first idea is to keep every edge $(i,j)$ with score above a threshold and run a weighted regression on the resulting edge-level file. Let
\[
w_{ij}=\frac{s_{ij}}{\sum_{j':(i,j')\in E}s_{ij'}}
\qquad\text{with}\qquad \sum_j w_{ij}=1.
\]
Define the score-weighted average covariate
\[
\tilde X_i=\sum_j w_{ij}X_j
\]
and the within-candidate-set variance
\[
V_i=\sum_j w_{ij}(X_j-\tilde X_i)^2.
\]

\begin{proposition}[Expanded versus collapsed slope]
Assume the regression is fit with centered variables and edge weights $w_{ij}$. The weighted least-squares slope from the expanded edge-level file is
\[
\hat\beta_{\mathrm{exp}}
=
\frac{\sum_i Y_i\tilde X_i}{\sum_i \tilde X_i^2+\sum_i V_i}.
\]
If instead one collapses to one row per source unit using covariate $\tilde X_i$, the resulting slope is
\[
\hat\beta_{\mathrm{col}}
=
\frac{\sum_i Y_i\tilde X_i}{\sum_i \tilde X_i^2}.
\]
Hence, whenever the two slopes share the same sign,
\[
|\hat\beta_{\mathrm{exp}}|\le |\hat\beta_{\mathrm{col}}|.
\]
\end{proposition}

\begin{proof}
Because $Y_i$ is constant across the candidate edges for source unit $i$,
\[
\sum_{i,j} w_{ij}Y_iX_j=\sum_i Y_i\tilde X_i.
\]
Also,
\[
\sum_{i,j} w_{ij}X_j^2=\sum_i \tilde X_i^2+\sum_i V_i
\]
by the usual variance decomposition identity.
\end{proof}

This is an algebraic result, not a subtle asymptotic one. Expanding the graph and reusing the same outcome across multiple candidate covariates mechanically adds the penalty $\sum_i V_i$ to the denominator. The more heterogeneous the candidate set, the larger the penalty.

The expanded file also corrupts the cluster structure. A target cluster in the expanded file contains outcomes generated under other targets' shocks whenever those other units include the target as a candidate. That makes the expanded graph a useful shipped artifact but a poor regression dataset.

\subsection{Score-collapsed regression}

A second natural idea is to collapse each candidate set to a single averaged covariate,
\[
\tilde X_i=\sum_j w_{ij}X_j.
\]
Write
\[
\tilde X_i = X_i^* + u_i,\qquad X_i^*\equiv X_{v(i)}.
\]
The quantity $u_i$ is measurement error induced by score averaging.

\begin{proposition}[Collapsed regression as nonclassical EIV]
Suppose
\[
E[\eta_{v(i)}+\varepsilon_i\mid X_i^*,u_i]=0.
\]
Then the probability limit of the collapsed OLS slope is
\[
\operatorname{plim}\hat\beta_{\mathrm{col}}
=
\beta_1
\frac{\operatorname{Var}(X^*)+\operatorname{Cov}(X^*,u)}
{\operatorname{Var}(X^*)+\operatorname{Var}(u)+2\operatorname{Cov}(X^*,u)}.
\]
If $\operatorname{Cov}(X^*,u)=0$, this reduces to classical attenuation. If $\operatorname{Cov}(X^*,u)\neq 0$, the bias can have either sign.
\end{proposition}

\begin{proof}
Write $Y_i=\beta_1 X_i^*+\nu_i$, where $\nu_i=\eta_{v(i)}+\varepsilon_i$ and $E[\nu_i\mid X_i^*,u_i]=0$. Then
\[
\operatorname{plim}\hat\beta_{\mathrm{col}}
=
\frac{\operatorname{Cov}(X^*+u,\beta_1X^*+\nu)}
{\operatorname{Var}(X^*+u)}
=
\beta_1
\frac{\operatorname{Var}(X^*)+\operatorname{Cov}(X^*,u)}
{\operatorname{Var}(X^*)+\operatorname{Var}(u)+2\operatorname{Cov}(X^*,u)}.
\]
\end{proof}

The intuition is straightforward. In many fuzzy matching problems, errors are not exchangeable across all records. They are local to a confusable neighborhood. A unit from \emph{rampur} is likely to be confused with \emph{rampura} or \emph{ramnagar}, not with an arbitrary distant place. Once ambiguity is structured in that way, $u_i$ will generally covary with $X_i^*$, and attenuation is no longer the default case.

\subsection{Magnitudes}

Table~\ref{tab:estimators} shows the five estimators in a medium-scale simulation ($K=40$ targets, $n=600$ source units, three confusable name clusters with covariate spread up to 0.50, DGP $\beta_1=5$). Table~\ref{tab:mi_diag} reports the MI diagnostics.

\begin{table}[ht]
\centering
\caption{Estimator comparison, medium-scale simulation ($K=40$, $n=600$).}
\label{tab:estimators}
\begin{tabular}{lrrrr}
\toprule
Method & $\hat\beta_1$ & SE & $N$ (rows) & Clusters \\
\midrule
True DGP & 5.000 & --- & --- & --- \\
Oracle & 6.322 & 1.233 & 600 & 40 \\
Canonical & 6.051 & 1.170 & 576 & 40 \\
Expanded & 3.655 & 0.329 & 1892 & 598 \\
Collapsed & 9.214 & 1.564 & 598 & 40 \\
MI (Rubin) & 6.005 & 1.151 & --- & --- \\
\bottomrule
\end{tabular}

\end{table}

\begin{table}[ht]
\centering
\caption{MI diagnostics via Rubin's rules ($M=200$ draws).}
\label{tab:mi_diag}
\begin{tabular}{lr}
\toprule
Quantity & Value \\
\midrule
Imputations ($M$) & 200 \\
Pooled $\hat\beta_1$ & 6.0054 \\
Within-imputation var ($\bar U$) & 1.3162 \\
Between-imputation var ($B$) & 0.0078 \\
Total var ($T$) & 1.3240 \\
Total SE & 1.1507 \\
$\lambda$ & 0.0060 \\
95\% CI & [3.750, 8.261] \\
\bottomrule
\end{tabular}

\end{table}

The collapsed estimator overshoots the oracle by 46\%. The expanded estimator undershoots by 42\%. MI recovers the oracle to within 5\%. The between-imputation share of variance is $\lambda = 0.006$, reflecting a design in which only 4\% of source units are ambiguous; in regimes with 15--35\% ambiguous records, $\lambda$ can exceed 0.10.

\section{Corrections}

\subsection{Multiple imputation over candidate assignments}

When the analyst has access to the candidate graph and calibrated probabilities, a direct fix is to treat the unresolved linkage as missing data. The required conditions are stronger than they are sometimes presented.

\begin{assumption}[Conditions for MI over candidate assignments]
The multiple-imputation strategy requires:
\begin{enumerate}
\item high recall: for most units, the true target lies in the candidate set;
\item calibrated probabilities: $p_{ij}$ is a reasonable approximation to $P(v(i)=j\mid \mathcal C_i)$;
\item an acceptable approximation to the dependence structure induced by the graph.
\end{enumerate}
\end{assumption}

Under these conditions, one can draw $M$ completed datasets by assigning one target per source unit in each draw, estimate the regression in each completed dataset, and combine results with Rubin's rules \citep{rubin1987}. This avoids both the within-candidate-set variance penalty of the expanded estimator and the averaging pathology of the collapsed estimator.

The diagnostic $\lambda = (1+1/M)B/T$, where $B$ is the between-imputation variance and $T$ is the total variance, is itself estimand-specific. If all ambiguous candidates fall within the same group or have similar covariate values, different draws produce similar estimates and $\lambda$ is small regardless of the ambiguity share. If the candidates straddle an estimand-relevant boundary, $\lambda$ concentrates exactly where it should. The same candidate graph produces different $\lambda$ values for different downstream analyses.

The main caveat is dependence. If a connected component contains many source units and a small number of targets, independent draws at the unit level may produce implausible assignments. In those settings, a component-level sampler or a joint posterior over matchings is preferable \citep{tancredi2011,steorts2016}.

\subsection{Validation-sample calibration}

If a validation sample with verified links is available, the problem changes. Let $X_i^\pi$ denote the covariate attached by the linked file and let $X_i^*$ denote the verified covariate in the validation sample. In the single-regressor linear case, suppose the calibration model is
\[
X_i^*=\alpha+\gamma X_i^\pi+r_i,\qquad E[r_i\mid X_i^\pi]=0.
\]
If the substantive model is linear with one mismeasured regressor, then a regression-calibration correction is
\[
\hat\beta_{\mathrm{RC}}=\frac{\hat\beta_{\mathrm{naive}}}{\hat\gamma}.
\]
This is a scalar correction for a scalar problem. With multiple regressors, interactions, or nonlinear outcome models, one needs multivariate regression calibration or a fuller joint error model rather than simple division \citep{carroll2006}.

The validation sample must represent the full estimand population, not just the subset retained by a canonical matching rule. Otherwise it cannot reveal selection induced by missed matches.

\section{Upstream contamination in shipped covariates}

Everything so far concerns the analyst's own linkage uncertainty. A separate issue arises when the target-table covariates were themselves created upstream by fuzzy matching.

Let $X_k^{\mathrm{pub}}$ denote the covariate published for target $k$. For a simple average, a useful stylized decomposition is
\begin{equation}
X_k^{\mathrm{pub}}
\approx
(1-\alpha_k)X_k^*+\alpha_k\bar X_{-k},
\label{eq:upstream}
\end{equation}
where $\alpha_k$ is the share of records assigned to $k$ that truly belong elsewhere, $X_k^*$ is the true target-level covariate, and $\bar X_{-k}$ is the mean among the wrongly assigned records. Equation \eqref{eq:upstream} is exact only when the correctly retained records for target $k$ are representative of the target; otherwise it should be read as a first-order approximation.

If the analyst then misassigns source unit $i$ to target $\pi(i)$ downstream, the total covariate error decomposes as
\[
u_i
=
\underbrace{(X_{\pi(i)}^*-X_{v(i)}^*)}_{\text{downstream wrong-target error}}
+
\underbrace{(X_{\pi(i)}^{\mathrm{pub}}-X_{\pi(i)}^*)}_{\text{upstream contamination}}.
\]
The analyst may be able to study the first term from the candidate graph they observe. The second term is frozen into the shipped product. Multiple imputation over the downstream graph therefore need not capture total linkage uncertainty.

This contamination is also estimand-specific. The same upstream matching error can severely distort one published covariate and barely move another, depending on within-neighborhood heterogeneity.

\section{Implications for data products}

A useful linked-data product should ship more than a final canonical linkage.

First, it should identify linkage status at the row level: exact match, unambiguous fuzzy match, ambiguous candidate set, or no candidate. Second, whenever feasible, it should expose the candidate graph or a compressed version of it. Without that object, analysts cannot perform estimand-specific re-resolution, multiple imputation, or meaningful sensitivity analysis. Because different estimands face different error costs (Section~3), different analysts may need to resolve the same candidate graph at different thresholds; shipping a single canonical linkage forecloses that choice. Third, it should document which covariates were themselves produced from earlier uncertain joins.

The broader lesson is that linkage uncertainty is not fully summarized by an overall match rate. What matters is where ambiguity lies, how it covaries with the variables used in the downstream estimand, and how much error was already embedded upstream in the shipped covariates.

\bibliographystyle{apalike}
\bibliography{validated_references}

\end{document}
